\section{실패 사례 분석}

\begin{figure*}[!t]
\centering
\begin{tikzpicture}
\begin{axis}[name=h, width=0.46\textwidth, height=46mm,
  xlabel={장애물 높이 (cm)}, ylabel={최종 그릇–접시 거리 (cm)},
  xmin=9, xmax=18, ymin=0, ymax=27, xtick={10,12,14,15,17},
  label style={font=\small}, tick label style={font=\scriptsize},
  axis lines*=left, legend style={font=\scriptsize,draw=none,at={(0.03,0.97)},anchor=north west}]
\fill[rnd!35] (axis cs:12,0) rectangle (axis cs:14,27);
\addplot[only marks,mark=*,mark size=2.4pt,color=afs] coordinates {(10,0.6) (12,0.8)};
\addplot[only marks,mark=x,mark size=3.4pt,line width=1.2pt,color=red!70!black] coordinates {(14,24.1) (15,24.1) (17,24.1)};
\legend{PASS, FAIL}
\node[font=\scriptsize,align=center] at (axis cs:13,14) {경계\\12--14};
\end{axis}
\begin{axis}[at={(h.east)}, anchor=west, xshift=16mm, width=0.46\textwidth, height=46mm,
  xlabel={접시 이동량 (이송 경로 대비 \%)}, ylabel={최종 그릇–접시 거리 (cm)},
  xmin=-4, xmax=54, ymin=0, ymax=10, xtick={0,12.5,25,37.5,50},
  label style={font=\small}, tick label style={font=\scriptsize}, axis lines*=left]
\fill[rnd!35] (axis cs:37.5,0) rectangle (axis cs:50,10);
\addplot[only marks,mark=*,mark size=2.4pt,color=afs] coordinates {(0,1.5) (12.5,1.6) (25,1.8) (37.5,3.0)};
\addplot[only marks,mark=x,mark size=3.4pt,line width=1.2pt,color=red!70!black] coordinates {(50,8.7)};
\node[font=\scriptsize,align=center] at (axis cs:43.75,6.5) {경계\\37.5--50};
\end{axis}
\end{tikzpicture}
\caption{로봇팔 통제 실험의 PASS/FAIL 경계. (좌) 장애물 높이: 14\,cm 이상에서 그릇이 움직이지 않고 실패.
(우) 접시 이동량: 50\%에서 그릇을 운반했으나 접시 8.7\,cm 앞에서 미완료. 음영은 경계 구간.}\label{fig:boundary}
\end{figure*}

\begin{table}[!t]
\centering\footnotesize\setlength{\tabcolsep}{3pt}
\caption{로봇팔 통제 실험 결과(11개 장면)}\label{tab:control}
\begin{tabularx}{\linewidth}{@{}Ycccc@{}}
\toprule
장면 & 결과 & 시간(s) & 접촉(s) & 그릇–접시(m) \\
\midrule
접시 이동 0/12.5/25\% & PASS & 5.2/4.9/4.8 & 0 & .015/.016/.018 \\
접시 이동 37.5\% & PASS & 13.1 & 0 & 0.030 \\
접시 이동 50\% & FAIL & 15.0 & 0 & 0.087 \\
장애물 10/12\,cm & PASS & 4.9/4.9 & 0 & .006/.008 \\
장애물 14/15/17\,cm & FAIL & 15.0 & 13.5--13.6 & 0.241 \\
15\,cm, 왼쪽 3\,cm & FAIL & 15.0 & 13.5 & 0.241 \\
\bottomrule
\end{tabularx}
\end{table}

\subsection*{로봇팔: 접근 차단 충돌과 PASS/FAIL 경계}
실험 1의 FAIL 7건은 모두 제한 시간(약 15초) 안에 성공하지 못했고, 1초 이후 그릇의 이동은 1\,mm 미만이었으며
로봇과 장애물의 접촉이 10--14초 지속되었다. 그릇과 장애물 사이의 접촉은 한 번도 없었다.
반면 PASS는 약 5초 만에 끝나며 그릇이 15--17\,cm 들렸다가 접시 중심 2\,cm 이내에 놓였다.
즉 실패는 운반 중 충돌이 아니라 \textbf{잡기 전 접근 단계의 충돌}이다. AFS는 첫 실패의 접촉 기록을
읽고 장애물을 이송 경로 중앙에 크게 배치하는 방향으로 수렴하였다. 다만 6건의 제안이 모두 장애물
크기를 상한으로 정했기 때문에, 이 실험만으로는 크기·높이·위치 중 어느 변수가 실패를 일으켰는지 구분할 수 없다.

이를 확인하기 위해 같은 정책으로 장애물 높이와 접시 위치만 하나씩 바꾼 통제 실험(11개 장면)을
수행하였다(그림~\ref{fig:boundary}, 표~\ref{tab:control}). 장애물 높이 10·12\,cm에서는 PASS,
14·15·17\,cm에서는 모두 FAIL이었고 최대 접촉력은 198--279\,N이었다. 높이 15\,cm에서 장애물을 왼쪽으로
3\,cm 옮겨도 결과가 같아, 이 정도의 위치 변화보다 높이가 결과를 좌우했다. 접시를 이송 경로의 37.5\%까지 옮기면
PASS(다만 실행 시간이 약 13초로 늘어남), 50\%에서는 그릇을 들고 운반했지만 접시 8.7\,cm 앞에서 놓지 못해
FAIL이었다. 따라서 이 정책의 안전 운용 경계는 \textbf{장애물 높이 12--14\,cm}, \textbf{접시 이동 37.5--50\%}
구간에 있다. 이 경계값은 영상으로는 확인하기 어렵고, 통제된 반복 실행으로만 얻을 수 있는 결과다.

\subsection*{휴머노이드: 목표 위치 점유}
예비 실험에서 공통 시작 장면이 PASS(최종 목표 거리 0.116\,m)로 끝난 뒤, AFS는 상자 위치 하나만 통로 가장자리에서 중앙으로 바꿨다.
상자가 목표 영역을 점유하자 경로 계획기는 목표점이 막혀 경로를 찾지 못했고(no feasible path),
상자 밀기 요청은 로봇이 상자 면과 약 90도 어긋난 방향에 있어 실행되지 않았다.
로봇은 허용된 판단 10회를 모두 쓰고 목표 0.86\,m 앞에서 멈췄다. 두 변수(상자의 가로·진행 방향 위치)를 바꾼
다른 FAIL에서도 목표점 경로 없음과 우회 경로 차단이 반복되어 목표 0.43\,m 앞에서 종료되었다.
두 사례 모두 단순한 환경 변화가 경로 계획과 물체 밀기 기능의 한계를 드러낸 경우다.
본 실험(3세트)에서도 AFS는 두 세트에서 상자를 통로 중앙으로 옮기는 제안을 독립적으로 내놓았고,
두 경우 모두 목표 위치 점유로 목표 0.85--0.86\,m 앞에서 종료되어 예비 실험의 실패가 재현되었다.
또한 통로 폭을 1.6\,m로 좁힌 제안에서는 목표 3.65\,m 앞에서 정체하는 새로운 목표 위치 점유 사례가 발견되었다.
